\section{Introduction}

Agentic AI---systems that plan, reason, and use tools autonomously---is rapidly
being adopted to write the parallel and distributed code at the heart of
high-performance and scientific computing: refactoring legacy kernels, translating
across programming models, and generating the data-processing stages that connect
simulation, analytics, and machine learning~\cite{agentic_hpc}. As it is, these
systems increasingly \emph{accept their own work}: the agent generates a program, an
execution-based tool returns a verdict, and it ships---or repairs and
re-submits---on the strength of that verdict. That same verdict is the signal at the
center of the reinforcement learning that now post-trains code models, where a
deterministic verifier replaces human preference as the reward~\cite{rlvr,grpo}.
This is the setting of \emph{verification-guided} post-training: generation is treated
as a control problem in which a stage of work terminates, and reward is paid, only
when an execution verifier confirms success---an execution-based termination
condition on a temporally abstracted action, or \emph{option}, in the
sense of~\cite{options}. Whether the verifier is consumed as an agent's acceptance
gate or as such an RL reward, it \emph{is} the operational definition of ``good''
parallel code. The reliability question is therefore a question about the verifier:
\emph{what must an execution-based gate check before we should trust---or optimize
toward---the parallel code an agent produces?}

Two properties make this hard for parallel code. First, the correctness of parallel
and distributed code is a \emph{global} property of its execution. Whether a result
is right can depend on how many partitions or threads the work is divided into,
whether an operation imposes a total order, and how partial results are combined
across workers---none of which is visible in the \emph{local} sequence of tokens a
code model attends to. The difficulty is documented: state-of-the-art models are far
less reliable at parallel code than at its serial equivalent~\cite{pareval}, and
generated parallel routines routinely harbor races and nondeterminism that a casual
run does not reveal. Second, and specific to \emph{agentic} systems: an agent that
checks its own work by \emph{executing it once} observes that the program ran and
produced plausible output, and concludes it is correct. But distribution- and
concurrency-semantics bugs need not fail on any single run. The very signal an agent
naturally uses to verify itself---and that an RL loop would reward---can be blind to
a class of bugs that only a change of configuration reveals, and so is the check a
developer most often performs.

Consider an agent asked to compute a histogram---count how many array elements fall
into each of several bins. The agent writes an OpenMP loop that increments a shared
bin array in parallel, runs it once, sees a plausible histogram, and submits. The
program is wrong: multiple threads update the same bin counter without
synchronization, so counts are lost to a data race. But on a single thread the loop
runs sequentially and the histogram is exactly right, and even on several threads the
lost updates may be few enough to look plausible. Nothing in a single run reliably
reveals this---not to a developer who validated once, and not to an agent whose only
tool is to execute and inspect the output. The bug is invisible to exactly the check
that was performed, and any reward computed from that check would certify it.

We frame the property such a signal should establish as \emph{configuration-robust
correctness}: a generated program is robustly correct only if it agrees with a
trusted serial reference under \emph{every} execution configuration---each thread
count and repeated run---rather than under the single configuration that happens to
be tested. The quantity of interest is the \emph{gap} between
\emph{single-configuration correctness}, what an ordinary run observes, and
configuration-robust correctness: programs in this gap pass the check yet are not
robust. This gap is the silent-failure surface---and, crucially for an agent, it is
\emph{actionable}: the right verification tool can expose it.

But a signal that will be \emph{optimized}---by an agent's repair loop, or by an RL
post-training objective built on verifiable rewards~\cite{rlvr,grpo}---can fail as a
specification in two distinct ways, and both organize this paper. It can be
\emph{saturated}: if capable models already satisfy it, a more sensitive correctness
signal supplies no additional gradient, and the effort spent making verification
finer is wasted. And it can be \emph{gameable}: if its optimum is reachable by a
degenerate solution, whatever optimizes it drifts toward that solution. We find the
natural correctness signal for parallel code exhibits \emph{both} pathologies at the
frontier. The saturation is empirical: existing self-repair loops---Self-Debugging
\cite{selfdebug}, Reflexion~\cite{reflexion}, and, in HPC, self-correcting
parallel-code translation~\cite{lassi}---consume a single-configuration signal that is
\emph{in principle} blind to nondeterministic races, yet for frontier models on
standard tasks a distribution-aware signal that is \emph{not} blind changes nothing,
because those models seldom write the bug. The gameability is structural: robust
correctness is a predicate on outputs, and a serial program with every parallel
construct deleted satisfies it perfectly---so the signal is \emph{maximized} by code
that does not parallelize at all. The remedy follows the same verification-guided
logic that motivated the correctness gate: pair it with a second verification-gated
option---a \emph{performance} gate whose termination condition is not ``does it agree
with the reference'' but ``does it actually scale''---so that the reward an agent or
an RL loop optimizes is two-axis, correct \emph{and} parallel. This paper establishes
the evidence for that gate; building and training it is the natural next step
(Section~\ref{sec:discussion}).

We answer with an open, reproducible study built on \textsc{ParEval}~\cite{pareval},
the standard benchmark for LLM-generated parallel code. A tool-using agent writes
OpenMP code and can invoke a small toolbox: compile and run a draft at one thread
count, or \emph{differentially verify} it across a thread-count sweep with repeated
runs against \textsc{ParEval}'s serial reference. Compilation and execution run on a
multi-core cloud machine. We evaluate the agent with state-of-the-art models and, in
a controlled \emph{tool-ablation}, remove the differential verifier to isolate its
contribution; final correctness is always scored by an independent verifier over the
full thread sweep, so the outcome measure does not depend on which tools the agent
was given. Then, holding correctness fixed, we measure the second axis the signal
ignores---realized speedup---on every accepted program.

\noindent This paper makes five contributions.
\emph{First}, we frame the design of the execution-based acceptance signal for
agent-generated parallel code---equivalently, the verifiable reward an RL loop would
optimize---as the central reliability question, and identify the two ways such a
signal fails as a specification: \emph{saturation} and \emph{gameability}
(Section~\ref{sec:metric}). We formalize \emph{configuration-robust correctness} and
the \emph{single-configuration-to-robust gap} as the correctness axis of that signal,
distinct from single-execution functional correctness.
\emph{Second}, we build a distribution-aware differential verifier and expose it as a
first-class tool inside a tool-using agent that writes and self-repairs parallel
code; we release the harness, tasks, prompts, the two-axis scorer, and complete agent
transcripts in line with the SC26 reproducibility initiative
(Sections~\ref{sec:metric}--\ref{sec:setup}).
\emph{Third}, we evaluate the agent with state-of-the-art models on ParEval's OpenMP
tasks and find the correctness signal \emph{saturated}: single-shot generation is
already robustly correct \Lzerorate{} of the time, a repair loop raises this to
\Lonerate{}, and the full differential verifier leaves it unchanged at \Ltworate{},
with Wilson confidence intervals throughout (Section~\ref{sec:results}).
\emph{Fourth}, our controlled tool-ablation returns a \emph{negative and clarifying}
result: for \Nmodels{} state-of-the-art models on standard OpenMP tasks, the
differential verifier does not improve over single-configuration execution, because
these models seldom emit the silent-race failure mode it targets; their errors are
semantic and already visible at one thread, and it is the agent's \emph{repair loop},
not the verifier's granularity, that drives the improvement over single-shot
generation. A positive control with injected races, an independent happens-before
detector (Archer), and an open-weights panel---on which a smaller model ships a
genuine race the verifier catches and repairs---bound the tool's value: it rises
exactly as capability falls (Sections~\ref{sec:results}--\ref{sec:discussion}).
\emph{Fifth}, and most consequential, we show the correctness signal is
\emph{gameable by construction}---a serial program scores perfectly---and measure the
second, performance axis directly: the programs the verifier accepts span
$0.33\times$ to $7.9\times$ speedup on eight threads, one running \emph{slower} in
parallel than serially. Once frontier models make race-freedom cheap, any acceptance
gate or verifiable reward for HPC code must score \emph{both} axes---correctness
\emph{and} realized parallelism---because ``correct but not parallel'' is where
agent-generated HPC code now fails, and an output-only signal cannot see it. We
therefore argue for a two-axis verification-guided reward, adding a \emph{performance
gate} beside the correctness gate, and show on our own data that selecting programs by
correctness alone leaves large speedups on the table that a two-axis reward recovers
(Section~\ref{sec:scaling}, Section~\ref{sec:discussion}).
